\section{Averaging the forests and separating the regular gas}
\label{sec:forest-average-log}
The exact block gas of Section~\ref{sec:gaussian-reassembly} was
constructed with a fixed forest. To restore the lattice symmetries we
average over forests before taking the logarithm of the flag-free gas.
We first show that this operation leaves the scalar normalization
unchanged, and then choose an average compatible with local factorization.

\subsection{The flat determinant is independent of the forest}
Every block forest has one outgoing edge from each nonroot vertex,
directed towards the root. Distinct nonroot vertices give distinct
unoriented positive-bond representatives of these edges. If two
outgoing edges represented the same bond with opposite orientations,
the forest height would strictly decrease in both directions, which
is impossible. Thus the tree bonds are in bijection with the nonroot
vertices. In particular all forests have the same number of tree
bonds and the same number $N$ of free bonds.

Let $F,H$ be two forests and choose a bijection $t$ between their
tree bonds. We construct an integer change of basis on all bond
coordinates, including the dummy tree coordinates. For a bond
cochain $u$, let $D_Fu$ set its $F$-tree coordinates to zero. Let
$\alpha_Fu(x)$ be the signed sum of $u$ along the forest path from
$x$ to its block root, and put
\[
                 B_Fu=u+\nabla\alpha_Fu.
\]
Telescoping along a path proves the following four facts:
$B_Fu$ is zero on the $F$-tree bonds; $B_Fu=u$ if $u$ was already
zero there; $\alpha_Fu$ is zero at each root; and
$B_F\nabla v=0$ when $v$ is zero at all roots. These assertions hold
over any abelian group, since they involve only signed sums.

Define $T_{FH}u$ on the $H$-tree bonds by transporting the old
$F$-tree coordinates through $t$, and on the other bonds by
$B_H D_Fu$. Then
\[
                    D_HT_{FH}u=B_HD_Fu.
\]
The reverse construction uses $t^{-1}$. On tree coordinates it
plainly recovers $u$. On the complementary coordinates,
\[
 B_FD_HT_{FH}u
   =B_F(D_Fu+\nabla\alpha_HD_Fu)=B_FD_Fu=D_Fu,
\]
because $\alpha_HD_Fu$ vanishes at the roots. Thus $T_{HF}T_{FH}=1$;
reversing the argument gives the other inverse identity. Its matrix
and inverse have integer entries. Their integer determinants multiply
to one, so $|\det T_{FH}|=1$. Apply this same construction to every
colour coordinate.

Let $M_F,M_H$ be the flat selected row matrices, and let $D_F^0,D_H^0$
be the diagonal projections onto their tree coordinates. Flat rows
annihilate gradients of root-zero functions. For a fine row this is
the telescoping identity around its plaquette, preserved by each
linear path coefficient. For a predictor row the coarse sampling
operator sends a gradient to the coarse gradient of its root values;
these values are zero. Signed coordinate changes preserve gradients,
so the statement holds simultaneously in all signed channels. Hence
\[
                 M_HT_{FH}=M_F,
        \qquad T_{FH}^TD_H^0T_{FH}=D_F^0.
\]
The second identity follows because the tree coordinates are simply
permuted. For the completed flat precisions we obtain
\[
 A_F^0=M_F^TM_F+D_F^0=T_{FH}^TA_H^0T_{FH},
                      \qquad \det A_F^0=\det A_H^0.
\]
Together with the common free-bond count, this proves that the scalar
$\mathscr O$ in \eqref{eq:reassembly-overhead} is forest independent.

\subsection{An independent choice at each coarse site}
There are finitely many allowed forests on a finite torus: each
vertex chooses one of finitely many oriented neighbouring edges and
one of finitely many heights. The set is nonempty by the block-tree
construction. At each coarse site $x$, choose independently a uniform
forest $F_x$ from this finite set. On the block with root $x$, use
only the restriction of $F_x$. The assembled object is a valid
forest, since each original forest edge remains in its parent block
and its height decreases along that edge. Repeated global forest
choices producing the same assembled forest cause no difficulty;
the average is defined on the product space of choices.

For a connected component $M$, its integrated activity depends on
these choices only at coarse sites in its full assigned support
$X_M$. This follows from the local reads of the rows, inverse
windows, tags and chart factors: agreement of the tree edges in
that support gives equality of each such read and of its Gaussian
marginal. To compare the Haar factors, realize both local charts as
pullbacks of the same full product Haar measure; agreement of the
principal Gaussian marginals then gives equality of the integrals.
In particular changing a forest outside $X_M$ does not
alter the component integral. Compatible components have disjoint
$X_M$. Independence of the site choices therefore gives
\begin{equation}\label{eq:forest-product-average}
 \mathbb E_F\prod_{M\in\mathcal A}a_M^F(V)
                    =\prod_{M\in\mathcal A}\mathbb E_Fa_M^F(V).
\end{equation}
This is the factorization of a product measure over disjoint
coordinate sets; it does not hold for an arbitrary average over
one global forest.

Set $\bar a_M(V)=\mathbb E_Fa_M^F(V)$. The bounds on $a_M^F$ are
uniform in $F$, so the same prices bound $\bar a_M$. The finite
average commutes with the absolutely convergent gas sum, and the
forest-independent $\mathscr O$ factors out. Thus
\eqref{eq:reassembly-block-gas} holds with $\bar a_M$ in place of
$a_M^F$. A signed lattice transformation permutes the site choices
and transforms every chosen forest by a bijection. Its product
uniform measure is unchanged. The covariance of the individual
sector formulas consequently gives the same covariance for
$\bar a_M$, with the original transformed activity addresses.

\subsection{The cost of excluding regular activities}
Separate the connected labels into regular labels $i$ with empty
inventory and forced labels $j$ with nonempty inventory. Write their
values as $w_i(V)$ and $k_j(V)$, and their supports as $X_i,Y_j$.
A nonzero regular value has no flag in its support. Fix a compatible
finite forced family $A$ whose inventories cover all flags, and put
$Y_A=\bigcup_{j\in A}Y_j$. The remaining sum is the regular gas
restricted to labels whose supports avoid $Y_A$:
\begin{equation}\label{eq:forest-forced-first}
 (\mathcal BR)(V)=\mathscr O(V)
       \sum_{A\ {\rm admissible}}\prod_{j\in A}k_j(V)\,Z_{Y_A}(V).
\end{equation}
The reindexing is exact: a compatible family splits uniquely into
its regular and forced parts, and cross compatibility means precisely
$X_i\cap Y_A=\varnothing$ for every retained regular label.
Absolute convergence justifies summing in this order.

Assume the regular prices satisfy the connected-cluster bound of
Section~\ref{sec:cluster-expansion}, with nonnegative support weight,
rooted allowance $\gamma$ and tail ratio $q_0<1$. Let $c_\Gamma$ be
the corresponding logarithmic term, indexed by a finite ordered
tuple of occurrences, with its Ursell coefficient divided by the
tuple factorial. Labels may repeat in this tuple. Then
\[
                    Z=\exp\left(\sum_\Gamma c_\Gamma\right)>0.
\]
For a finite set $Y$ define
\[
        P_Y=\sum_{\Gamma:\,X_\Gamma\cap Y\ne\varnothing}c_\Gamma,
               \qquad X_\Gamma=\bigcup_{i\in\Gamma}X_i.
\]
Deleting labels that meet $Y$ replaces $c_\Gamma$ by zero precisely
when its union support meets $Y$. Otherwise it leaves the term
unchanged. The same absolute majorant applies to the deleted gas.
Subtracting the two absolutely convergent logarithmic series yields
\begin{equation}\label{eq:forest-deletion-ratio}
                         Z_Y=Z e^{-P_Y}.
\end{equation}
The rooted cluster estimate, summed over $x\in Y$, gives
\begin{equation}\label{eq:forest-deletion-price}
 \sum_{\substack{\Gamma:\,X_\Gamma\cap Y\ne\varnothing\\
                       |\Gamma|-1\ge m}}|c_\Gamma|
          \le |Y|\frac{\gamma q_0^m}{1-q_0}.
\end{equation}
Indeed each term on the left is counted by at least one site of
$Y\cap X_\Gamma$; overcounting is harmless in this nonnegative sum.
In particular, with $C_Y=|Y|\gamma/(1-q_0)$,
\[
                Ze^{-C_Y}\le Z_Y\le Ze^{C_Y}.
\]
The correction $P_Y$ retains its actual sign; only its absolute
estimate uses $C_Y$. Combining \eqref{eq:forest-forced-first} and
\eqref{eq:forest-deletion-ratio} gives the exact separation
\begin{equation}\label{eq:forest-regular-separated}
 (\mathcal BR)(V)=\mathscr O(V)Z(V)
       \sum_{A\ {\rm admissible}}\prod_{j\in A}k_j(V)e^{-P_{Y_A}(V)}.
\end{equation}
For each fixed $A$, the ratio is finite and positive. The original
absolute family sum implies absolute summability of this outer
series: its term is exactly the signed forced product times the
restricted regular gas, whose absolute value is bounded by the
corresponding sum of absolute original terms. No volume-independent
bound on the extensive scalar $Z$ is asserted.

\subsection{Local regular functions and their derivatives}
For each regular component, the integrated construction gives a
measurable local extension $f_i$, smooth through the purchased finite
order at configurations healthy on $X_i$, with
\[
w_i(V)=\mathbf1_{\{X_i\cap\operatorname{Flags}(V)=\varnothing\}}f_i(V).
\]
The forest average preserves locality, gauge invariance, measurability
and each uniform derivative bound, by linearity and the triangle
inequality for the finite average. For a cluster tuple
$\Gamma=(i_1,\ldots,i_n)$, define
\[
 F_\Gamma(V)=\frac{u(i_1,\ldots,i_n)}{n!}
                                      \prod_{r=1}^n f_{i_r}(V).
\]
Then
\begin{equation}\label{eq:forest-cluster-mask}
 c_\Gamma(V)=
 \mathbf1_{\{X_\Gamma\cap\operatorname{Flags}(V)=\varnothing\}}
                                                    F_\Gamma(V).
\end{equation}
If the union is healthy, every constituent mask is one. If it is
not, a flagged site belongs to some constituent support, so that
factor vanishes. This proves the equality even for repeated labels.

Suppose the local extension of the $r$th occurrence has derivatives
through order $j$ bounded by $p_{i_r}$ at derivative scale one at
configurations healthy on its support. Restrict throughout the
following derivative estimates to configurations healthy on
$X_\Gamma$; every constituent hypothesis then holds.
Leibniz's formula gives product amplitude $\prod_rp_{i_r}$ at
derivative scale $n$. Thus its derivatives through order $j$ are
bounded by
\[
        \frac{|u(i_1,\ldots,i_n)|}{n!}
                    (1+n)^j\prod_rp_{i_r}
       \le\frac{|u(i_1,\ldots,i_n)|}{n!}
                                      \prod_r(2^jp_{i_r}),
\]
since $1+n\le2^n$ for $n\ge1$. Buying the regular cluster bound
with the inflated prices $2^jp_i$ therefore pays these derivatives
on the same union support. A spatial derivative weight is also local:
from any anchor in $X_\Gamma$, the connected-cover diameter bound
places every read within a fixed multiple of
$(B+1)(1+\operatorname{cover}_B X_\Gamma)$. Its polynomial cost is
paid by the reserved exponential support weight as in
Section~\ref{sec:mixed-activity-prices}.

Equations~\eqref{eq:forest-regular-separated} and
\eqref{eq:forest-cluster-mask} supply the regular logarithm and the
exact exclusion correction of the nonlinear step. The exclusion
exponential must still be regrouped with the forced labels using
its occurrence multiplicities; it cannot be dropped or replaced by
an independent probability. Its quantitative allowance is the
rooted price \eqref{eq:forest-deletion-price}.
